% Source: https://artofproblemsolving.com/community/c3h2542700_2005_national_team_round_2
% Stable identity: aops-topic-2542700-q01
Let $v$ be the number of votes and $c$ the number for the losing candidate. Rounding to $31\%$ means
\[30.5\%\le\dfrac cv<31.5\%,\qquad\text{i.e.}\qquad 0.305\,v\le c<0.315\,v.\]
The interval $[0.305v,\,0.315v)$ has length $0.01v<1$, so we can test $v=1,2,3,\dots$ in turn. For $v=1,2,\dots,12$ the intervals are
\[[0.305,0.315),\ [0.61,0.63),\ [0.915,0.945),\ [1.22,1.26),\ [1.525,1.575),\ [1.83,1.89),\]
\[[2.135,2.205),\ [2.44,2.52),\ [2.745,2.835),\ [3.05,3.15),\ [3.355,3.465),\ [3.66,3.78),\]
and none of them contains an integer. For $v=13$ it is $[3.965,4.095)$, which contains $c=4$: indeed $\dfrac4{13}\approx30.77\%$, which rounds to $31\%$, and $4<\dfrac{13}{2}$, so this candidate does lose.

The minimum number of votes is $\boxed{13}$.
